\documentclass[12pt,a4paper]{article}

\usepackage[cm]{fullpage}
\usepackage{amsthm}
\usepackage{amsmath}
\usepackage{amsfonts}
\usepackage{amssymb}
\usepackage{xspace}
\usepackage[english]{babel}
\usepackage{fancyhdr}
\usepackage{titling}
\usepackage{framed}
\usepackage{stmaryrd}
\usepackage{mathpartir}
\usepackage{stix}
\usepackage{xcolor}
%\renewcommand{\thesection}{Task \projnumber.\arabic{section}:}

%%%%%%%%%%%%%%%%%%%%%%%%%%%%%%%%%%%%%%%%%%
%         exercise-specific macros       %
%%%%%%%%%%%%%%%%%%%%%%%%%%%%%%%%%%%%%%%%%%

\newcommand{\mcall}[2]{#1(#2)}
\newcommand{\call}[2]{\mcall{\textit{#1}}{#2}}
\newcommand{\step}{\rightarrow}
\newcommand{\nstep}{\overset{*}{\rightarrow}}
\newcommand{\update}[3]{#1[#2 \mapsto #3]}
\newcommand{\run}[2]{\llbracket #1 \rrbracket #2}
\newcommand{\loweq}{\mathbin{=_L}}
\newcommand{\tainted}[1]{\textcolor{red}{#1}}
\newcommand{\untainted}[1]{\textcolor{blue}{#1}}

\newcommand{\config}[1]{\ensuremath{\left\langle #1 \right\rangle}\xspace}
\newcommand{\ttt}[1]{\texttt{#1}}
\newcommand{\proves}{\vdash}
\newcommand{\sth}{\ensuremath{.\;}\xspace}
\newcommand\set[1]{\ensuremath{\{#1\}}}
\newcommand{\WHILE}[1]{\ensuremath{\mathsf{while}\ #1 \ \mathsf{do}}\xspace}
\newcommand{\IF}{\ensuremath{\mathsf{if}}\xspace}
\newcommand{\THEN}{\ensuremath{\mathsf{then}}\xspace}
\newcommand{\ELSE}{\ensuremath{\mathsf{else}}\xspace}
\newcommand{\SKIP}{\ensuremath{\mathsf{skip}}\xspace}

%%%%%%%%%%%%%%%%%%%%%%%%%%%%%%%%%%%%%%%%%%
% This part needs customization from you %
%%%%%%%%%%%%%%%%%%%%%%%%%%%%%%%%%%%%%%%%%%

% please enter your name and matriculation number here
%DONE
\newcommand{\name}{Benjamin Deutsch}
\newcommand{\matriculation}{12215881}

%%%%%%%%%%%%%%%%%%%%%%%%%%%%%%%%%%%%%%%%%%
%           End of customization         %
%%%%%%%%%%%%%%%%%%%%%%%%%%%%%%%%%%%%%%%%%%

\newcommand{\projnumber}{3}
\newcommand{\Title}{IFC and Noninterference}
\setlength{\headheight}{15.2pt}
\setlength{\headsep}{20pt}
\setlength{\textheight}{680pt}
\pagestyle{fancy}
\fancyhf{} 
\fancyhead[L]{Project \projnumber\ -- \Title}
\fancyhead[C]{}
\fancyhead[R]{\name{} (\matriculation)}
\renewcommand{\headrulewidth}{0.4pt}
\fancyfoot[C]{\thepage}


\begin{document}
\thispagestyle{empty}
\noindent\framebox[\linewidth]{%
 \begin{minipage}{\linewidth}%
 \hspace*{5pt} \textbf{Formal Methods for Security and Privacy (SS26)} \hfill Prof.~Matteo Maffei \hspace*{5pt}\\

 \begin{center}
  {\bf\Large Project \projnumber~-- \Title}
 \end{center}

 \vspace*{5pt}\hspace*{5pt}\name{} (\matriculation) \hfill TU Wien \hspace*{5pt}
\end{minipage}%
}
\vspace{0.5cm}
%%%%%%%%%%%%%%%%%%%%%%%%%%%%%%%%%%%%%%%%%%%%%%%

% README
% remember to fill the macros \matriculation and \name


\setcounter{section}{3}
\section{Noninterference Analysis}
\section{Relaxing Noninterference}
\subsection{Tax Bracket Calculation}
Nia analysis of program in Figure 4 will output that the program is interferant.
As in the lecture, we can represent the states of the memory for each variable using
$\oplus$ $\circledcirc$ and $\ominus$. Consider the following two states of the memory:
\begin{align*}
    m_1 &: \{\{\untainted{lwinner: \oplus\circledcirc\ominus}, \tainted{hbidA: \ominus}, \tainted{hbidB: \oplus}\}\} \\
    m_2 &: \{\{\untainted{lwinner: \oplus\circledcirc\ominus}, \tainted{hbidA: \oplus}, \tainted{hbidB: \ominus}\}\}.
\end{align*}

To process line 1, we apply Propagate-Taint on both traces to get the following states:
\begin{align*}
    m_1 &: \{\{\untainted{lwinner: \ominus}, \tainted{hbidA: \ominus}, \tainted{hbidB: \oplus}\}\} \\
    m_2 &: \{\{\untainted{lwinner: \ominus}, \tainted{hbidA: \oplus}, \tainted{hbidB: \ominus}\}\}.
\end{align*}

To process line 3, we need to apply the Raise-Context-Step rule, since the control flow depends on
tainted variables. Therefore, until joining at line 11, the context is now tainted. Note that 
the trace starting with $m_1$, is now before line 4, since the condition is true in that trace. 
The other trace is before line 6, since the condition is false in that trace.

Now, we can apply Propagate-Taint on line 4 for the first trace to get 
\begin{align*}
    m_1 &: \{\{\tainted{lwinner: \oplus}, \tainted{hbidA: \ominus}, \tainted{hbidB: \oplus}\}\}. \\
\end{align*}

$lwinner$ is now tainted, since it was assigned under a tainted context. At this point, the first 
trace is before line 11, waiting to join with the second trace. 

On the second trace, we can apply Propagate-Taint to process line 6, since the current context is 
already tainted. As the condition is true, the trace is now before line 7. To process line 7, we can 
apply Propagate-Taint again to get
\begin{align*}
    m_2 &: \{\{\tainted{lwinner: \circledcirc}, \tainted{hbidA: \oplus}, \tainted{hbidB: \ominus}\}\}. \\
\end{align*}

Now, both traces are before line 11. Therefore, we can apply the Join rule, to get
\begin{align*}
    m_1 &: \{\{\tainted{lwinner: \oplus}, \tainted{hbidA: \ominus}, \tainted{hbidB: \oplus}\}\} \\
    m_2 &: \{\{\tainted{lwinner: \circledcirc}, \tainted{hbidA: \oplus}, \tainted{hbidB: \ominus}\}\}. \\
\end{align*}

Note that, since the value of $lwinner$ differs in the two traces and it is tainted in both traces, 
it remains tainted after the join. After processing the two remaining assignments using Propagate-Taint, 
we get the following states:
\begin{align*}
    m_1 &: \{\{\tainted{lwinner: \oplus}, \untainted{hbidA: \circledcirc}, \untainted{hbidB: \circledcirc}\}\} \\
    m_2 &: \{\{\tainted{lwinner: \circledcirc}, \untainted{hbidA: \circledcirc}, \untainted{hbidB: \circledcirc}\}\}. \\
\end{align*}

Since the public variable $lwinner$ is tainted, the program is interferant.

\subsection{Controlled Information Leaks: Abstract Noninterference}
My implementation first checks whether the variables $hbidA$ and $hbidB$ are present in the variables list.
If so, then $hbidA$ and $hbidB$ are added to the observed variable list, since those are needed to compute 
$\phi$. Also, if $hbidA$ and $hbidB$ are present, auction\_winner will returns a z3 expression calculating
the winner of the auction. If $hbidA$ and $hbidB$ are not present, the the function returns $0$. Finally,
also relate\_executions checks for the presence of $hbidA$ and $hbidB$. If they are not present, it simply 
checks low variable equivalence. If $hbidA$ and $hbidB$ are present, then it additionally checks that the 
winner of the auction is the same in both traces.

\paragraph{simple\_winner.while.} This program is wrongly labeled as interferent, since the if statement is 
the last statement of the program. The abstraction of $\mathrm{Seq}(c_1,c_2)$ only checks whether the $c_1$ 
might require a join, but not whether $c_2$ requires one. Therefore, if the statement requiring a join is 
the last statement of the program, the $\mathrm{Join}$ rule is never applied, which is why $lwinner$ stays 
tainted in the end.

To fix this, we could 
\begin{itemize}
    \item change the $\mathrm{Seq}$ abstraction to also check whether $c_2$ is a possible join point, or
    \item add a trivial statement (e.g. a self-assignment, like it is done in simple.while) after 
          the if statement, enabling the $\mathrm{Join}$ rule to be applied.
\end{itemize}

\end{document}

